\documentclass[sigplan,11pt,nonacm=true]{acmart}
\settopmatter{printfolios=false,printccs=false,printacmref=false}

\usepackage{booktabs}
\usepackage{threeparttable}

% v2: camera-ready version with re-run results.
\newcommand{\tool}{archcheck}
\newcommand{\x}{$\times$}

\begin{document}

\title{Does AI-Assisted Development Accelerate Architectural Drift?\\
A Large-Scale Study of C++ Repositories}

\author{Sergey Blekher}
\affiliation{%
  \institution{Independent Researcher}
  \country{Russia}
}
\email{sergey.blekher@gmail.com}

\begin{abstract}
Coding agents can generate software changes faster than human reviewers can
inspect them, raising the concern that repository structure degrades faster
under AI-assisted development. We test this hypothesis in C++ by measuring
architectural drift at commit level across 1{,}188 open-source repositories and
484{,}500 analyzed commits, using a deterministic source-level analyzer that
replays history and records structural regressions: new or grown include
cycles, god headers, cross-area dependencies, code clones,
cognitive-complexity growth, and self-admitted technical debt. AI attribution
relies only on repository metadata, never on commit prose.

A naive pooled comparison suggests that AI-attributed commits introduce
substantially more complexity. This result disappears after controlling for
repository and commit size: in a within-repository, size-matched analysis over
344 repositories, we detect no drift category in which AI-attributed commits
are worse than human commits. Attribution relies on visible markers, so the
human class contains unmarked AI use; under the contamination we model, a
true 2\x{} per-commit penalty would be observed as about 1.4\x{}, and the
sparse structural categories have low power. The observable
difference is commit size: AI-attributed commits are about 1.3\x{} larger.

In a preregistered adoption event study covering 854 adopting repositories and
517 placebo controls, agent adoption is associated with 1.93\x{} more drift
events per week than in placebo repositories observed in the same calendar
weeks (95\% CI $[1.42, 2.62]$). Per commit the increase is 1.24\x{}
($[0.97, 1.58]$) and per KLOC 1.14\x{} ($[0.96, 1.35]$); cognitive-complexity
growth is the one endpoint that also rises per unit of change (1.17\x{} per
commit, 1.14\x{} per KLOC). Merged pull requests rise 1.38\x{} relative to
placebo, while human review comments per pull request fall to 0.73\x{}.

Most of the post-adoption increase in drift comes from more changes arriving,
with a smaller per-change rise in complexity; the number of human review
comments does not grow with the number of pull requests.
\end{abstract}

\keywords{architectural drift, AI-assisted development, static analysis, C++,
mining software repositories, code review}

\maketitle

\section{Introduction}

The current debate about AI-assisted software development often assumes a
direct quality effect: if coding agents are less reliable than humans,
repositories that use them should accumulate more technical and architectural
debt per unit of change. This is plausible. Repository-wide constraints are
easy to omit from a local prompt, while coding agents are optimized to
complete the immediate task; controlled experiments show agent performance
degrading substantially as structural constraints accumulate
(\emph{constraint decay}~\cite{dente2026constraintdecay}). In C++, where
physical design is expressed partly through headers and include
dependencies~\cite{lakos1996}, a locally reasonable edit can create a
non-local regression such as an include cycle or a new dependency across
subsystems.

The opposite possibility is less dramatic but operationally important. Agents
may produce code of roughly the same structural quality per unit of change
while increasing throughput. If the amount of code and the number of changes
rise faster than human review capacity, then the absolute rate of
architectural drift can still increase even without an intrinsic quality
penalty in agent-authored code.

This paper distinguishes these two mechanisms. We study architectural drift at
three levels: individual commits, repository adoption events, and review
capacity. Our analysis covers 1{,}188 C++ open-source repositories and
484{,}500 analyzed commits. Rather than infer architectural quality from
repository snapshots alone, we replay history and measure the delta introduced
by each commit.

We ask three research questions:

\begin{itemize}
\item \textbf{RQ1.} Do AI-attributed commits introduce more architectural
  drift than human commits of comparable size within the same repository?
\item \textbf{RQ2.} Does the rate of architectural drift change after a
  repository adopts coding agents?
\item \textbf{RQ3.} Does recorded human review capacity scale with the
  post-adoption increase in development volume?
\end{itemize}

Our main result is a separation between quality per unit and change volume. In
size-matched within-repository comparisons, we do not observe an agent
disadvantage in any measured drift category, and agent pull requests are not
dirtier at their final submitted head (Section~\ref{sec:p5}). After adoption, however, drift events accumulate
faster per week, mostly because more changes arrive, and the number of human review
comments does not grow with the number of pull requests. The practical risk is therefore not simply that agents write
structurally worse code. It is that they increase the rate of structural
change faster than existing review processes absorb it.

The paper makes three contributions. First, it provides a large-scale
per-commit census of architectural drift in C++ repositories, including
physical-design regressions (include cycles, god headers, cross-area
dependencies) that repository-snapshot studies do not measure. Second, it
separates the author effect from the throughput effect using within-repository
size matching, an as-submitted pull-request analysis, and a preregistered
adoption event study with placebo controls, with a direction check on another
study's repository list. Third, it connects the
observed volume increase to review capacity, showing that merged pull requests
increase after adoption while human review comments per pull request decrease.

\section{Measurement Instrument}
\label{sec:instrument}

We built \tool{}\footnote{\url{https://github.com/blurman-ai/archcheck}, open
source. Every rule has passing and failing fixture projects in the repository's
test suite.} for this study: a deterministic static analyzer that measures
structural regressions between two repository revisions. It does not call an
LLM. The analyzer reads C and C++ source trees directly from Git objects,
constructs a physical dependency graph, computes local structural metrics, and
subtracts the pre-change state from the post-change state.

For include structure, every resolved include becomes a directed edge between
source files. Strongly connected components are detected with Tarjan's
algorithm~\cite{tarjan1972}; a newly created or enlarged non-trivial component
is recorded as cycle drift. Headers crossing a fan-in threshold of 50 are
recorded as god-header drift. Cross-area dependencies are derived from
repository paths and recorded when a change introduces a dependency between
previously separate areas. These checks operationalize the physical-design
rules of Lakos~\cite{lakos1996} and the source-file (SF) section of the C++
Core Guidelines~\cite{coreguidelines}.

Local code signals include cognitive-complexity
growth~\cite{campbell2018}, newly added TODO/FIXME/HACK markers
(self-admitted technical debt), lack of test co-evolution, and token-based
clone introduction. Clone detection normalizes identifiers and literals before
candidate matching, so copied-and-renamed fragments are still detected. The
analyzer excludes bulk imports and classifies vendored, generated, and test
code before applying authored-code rules.

The unit of analysis is a commit. For each eligible commit $c$ and drift
category $k$, we define $D(c,k)=1$ when the commit introduces at least one new
finding in category $k$. This design intentionally measures \emph{regression}
rather than repository debt stock: a repository may contain old cycles or
duplication without charging them to every subsequent commit.

\paragraph{Instrument validity.} Output is deterministic: byte-identical
across Git backends and across re-runs on the same input. Three measurement
bugs found during development (an empty-blob desynchronization in the
in-memory Git backend, shallow-clone boundary commits read as whole-tree
additions, inconsistent authored-scope filtering) were fixed, and the affected
analyses were rerun; on 128{,}393 commits in 226 repositories the corrected
replay kept the RQ1 verdict in six of seven categories, and the seventh moved
towards agents. Against NiCad on an independent C benchmark (monit, function
granularity) the clone detector covered 8 of 12 NiCad clone classes, and all
21 of its reported pairs were true clones on manual inspection. Per-commit
clone introductions had 86--91\% precision in a corpus audit. Per-category
manual audits are in the artifact.

\section{Data and Attribution}
\label{sec:data-attribution}

Repositories entered the corpus if they were written primarily in C or C++,
had more than 300 commits since 1~May~2025, showed some agent activity (an AI
\texttt{Co-authored-by} trailer, at least 10\% bot-authored commits, or a
median commit-message length of at least 100 characters), were not forks, and
were not large projects with enforced review (CODEOWNERS, branch protection,
many maintainers). About 2.7\% of the screened repositories passed. The corpus
is therefore enriched for agent adopters and excludes the projects with the
strictest review; it is not a random sample of C++ open source. Commits since
June~2024 were analyzed.

The primary corpus contains 1{,}188 C++ open-source repositories and 484{,}500
successfully analyzed commits out of 520{,}177 attempted. Of the commits not
analyzed, 33{,}582 belong to repositories that exceeded the time budget and
790 timed out individually, so coverage is weighted towards faster
repositories. Of these, 72{,}778
commits (15.0\%) are AI-attributed, 408{,}716 are human-attributed, and
3{,}006 are non-AI automation. AI attribution uses only observable metadata:
bot authorship, AI-associated committers, or explicit \texttt{Co-authored-by}
trailers. Commit message prose is never used as a classifier.

This attribution strategy is deliberately conservative. It has high precision
for visible agent use but incomplete recall; multi-method census work shows
that single-signal detection can undercount specific agents by more than an
order of magnitude~\cite{khosravani2026census}. Human-attributed commits
therefore include an unknown amount of unmarked AI assistance, such as IDE
completion or stripped trailers: ``human-attributed'' means ``no observable
AI marker''. This contamination biases author comparisons
toward the null and is quantified in Section~\ref{sec:threats}.

Commit size is the major confounder. In the corpus, the median AI-attributed
commit changes 50 lines versus 21 for human-attributed commits; within
repositories, the median AI/human size ratio is 1.26\x{} ($p=0.005$). Because
several drift signals grow mechanically with change size, pooled author
comparisons are not interpretable without controlling for repository and
commit size.

\section{RQ1: Agent Commits versus Human Commits}
\label{sec:rq1}

\subsection{Design}

We compare AI-attributed and human-attributed commits inside the same
repository, matching commits by size band on changed lines and aggregating
paired agent-minus-human differences at repository level. The analysis
includes the 344 repositories with sufficient observations from both classes
within the matched bands. We use a two-sided sign test over repository-level
differences; this treats each repository as one observation and prevents a
handful of very large repositories from dominating the result. With seven
category-level tests, the Bonferroni-adjusted significance threshold is
$\alpha = 0.05/7 \approx 0.0071$.

\subsection{Size control changes the answer}

Table~\ref{tab:ladder} shows why the control structure matters, using the
complexity category. Pooled, 27.0\% of AI-attributed commits trigger the
complexity-growth signal versus 15.4\% of human commits ($\approx$1.75\x).
A raw within-repository comparison shrinks this to $+1.55$ percentage points
($p=0.037$). After size matching, the median difference is $+0.06$ pp
($p=0.44$): the apparent complexity penalty is explained by larger commits,
not by a different per-change propensity to add complexity.

\begin{table}[t]
\caption{The complexity comparison at three levels of control. The pooled
``headline'' vanishes once repository and commit size are held fixed.}
\label{tab:ladder}
\footnotesize
\begin{tabular}{@{}lll@{}}
\toprule
Control level & Complexity result & $p$ \\
\midrule
Pooled (none) & 27.0\% vs.\ 15.4\% ($\approx$1.75\x) & --- \\
Within-repo, raw & $\Delta = +1.55$ pp (135/102 repos) & 0.037 \\
Within-repo, size-matched & $\Delta = +0.06$ pp (172/157 repos) & 0.44 \\
\bottomrule
\end{tabular}
\end{table}

\subsection{All seven categories}

Table~\ref{tab:seven} reports the same size-matched within-repository analysis
for every drift category. There is no category in which AI-attributed commits
are measurably worse after matching.

\begin{table*}[t]
\caption{Within-repository, size-matched comparison over 344 repositories.
$\Delta$ is the median repository-level agent-minus-human difference in the
share of commits with a finding (percentage points). ``A$>$H''/``A$<$H''
count repositories where agents trigger the category more/less often;
sign-test $p$ is two-sided over non-tied repositories. ``A$>$H bound'' is the
one-sided 95\% upper bound (Clopper--Pearson) on the share of non-tied
repositories where agents are worse; ``detectable'' is the smallest such share
the test detects with 80\% power at $\alpha=0.05/7$.}
\label{tab:seven}
\small
\begin{threeparttable}
\begin{tabular}{@{}lrrrrrrl@{}}
\toprule
Category & $\Delta$ (pp) & A$>$H / A$<$H & Ties & $p$ & A$>$H bound & detectable & Verdict \\
\midrule
complexity      & $+0.06$ & 172 / 157 & 15  & 0.44               & 0.57 & 0.60 & null \\
new clone       & $-0.23$ & 130 / 178 & 36  & 0.0073             & 0.47 & 0.60 & cleaner (nominal) \\
SATD            & $-0.52$ & 85 / 191  & 68  & $1.6\cdot10^{-10}$ & 0.36 & 0.61 & agents cleaner\tnote{\dag} \\
test co-evo     & $-1.63$ & 121 / 206 & 17  & $3\cdot10^{-6}$    & 0.42 & 0.60 & agents cleaner\tnote{\dag} \\
cross-area dep  & $+0.00$ & 60 / 104  & 180 & 0.0007             & 0.43 & 0.64 & no deficit, sparse \\
grown cycles    & $+0.00$ & 8 / 30    & 306 & 0.0005             & 0.35 & 0.78 & insufficient events \\
god headers     & $+0.00$ & 28 / 53   & 263 & 0.0073             & 0.44 & 0.70 & insufficient events \\
\bottomrule
\end{tabular}
\begin{tablenotes}\footnotesize
\item[\dag] Stylistic explanations are possible; see Section~\ref{sec:rq1}.
\end{tablenotes}
\end{threeparttable}
\end{table*}

Two clarifications are necessary to read the sparse structural rows correctly.
First, a median $\Delta$ of $0.00$ next to a significant $p$-value is not a
contradiction: the sign test is computed over the minority of repositories
that are \emph{not} tied, and among those the direction favors agents; the
median repository simply has no such event for either class. For cross-area
dependencies, cycles, and god headers, 180--306 of 344 repositories are ties.
We therefore read these rows as \emph{no agent deficit}, not as evidence that
agents are population-wide cleaner. Second, four results (SATD, test
co-evolution, cross-area dependencies, grown cycles) survive the
Bonferroni-adjusted threshold, while the copy-paste and god-header differences
($p = 0.0073$) are nominally significant but narrowly miss it
($\alpha \approx 0.0071$); the table labels them accordingly. For
perspective, the only nominally significant result \emph{against} agents
anywhere in this analysis --- the raw within-repository complexity difference
of Table~\ref{tab:ladder} ($p = 0.037$) --- misses the same threshold by a
far wider margin.

The two ``agents cleaner'' rows have plausible explanations unrelated to
design quality. Agents rarely write TODO/FIXME/HACK markers, which lowers
SATD by style alone, and agents are often instructed to update tests in the
same change, which raises test co-evolution. We do not read either row as
evidence that agent code is structurally better.

Two further checks corroborate the null on the structural categories. A
commit-level null model that attributes drift births within each repository
proportionally to each class's commit share finds AI commits
\emph{under}-represented among the commits that create them: 17 grown-cycle
births against an expectation of 25.9 (0.66\x), and 194 cross-area events
against an expectation of 427.9 (0.45\x). And the corrected-scanner replay of
Section~\ref{sec:instrument} reproduces the verdict of
Table~\ref{tab:seven} in six of seven categories, with the seventh moving in
favor of agents.

\subsection{As-submitted pull requests}
\label{sec:p5}

Every estimate above is measured on \emph{merged} history: only commits that
survived review. If reviewers caught agent drift before merge, agents would
look clean in merged history because a human gate stood in front of them. To
close this gap we measured drift on the state agents \emph{submitted}
(\texttt{pull/N/head}), comparing 291 agent-authored and 260 human-authored
merged pull requests within the same 18 repositories (head refs recovered for
99.8\% of PRs). The pooled rates look agent-dirtier (55.3\% vs.\ 45.4\% of PRs
carry at least one drift finding), but this is a repository-composition
artifact: agent PRs concentrate in a few high-drift repositories. The
registered within-repository contrast, the median of per-repository
agent-minus-human differences, is $-0.210$ with a repo-clustered bootstrap
CI of $[-0.365, 0.000]$ and a leave-one-repo-out range of $[-0.235, -0.185]$:
slightly on the ``agent cleaner'' side, with the CI touching zero. Two
limits apply. The sample contains merged pull requests only, so drift in
pull requests that were closed without merging is not observed. And \texttt{pull/N/head} is the final head of
the branch, which includes revisions made in response to review. The analysis
therefore shows that merged agent PRs are not dirtier at their final head;
it does not exclude review filtering.

\medskip
\noindent\textbf{RQ1 answer.} We find no evidence that AI-attributed commits
introduce more architectural drift than human commits of comparable size
within the same repository, and none that merged agent pull requests are
dirtier at their final head. The dominant observable difference is commit size, not drift
propensity per comparable change.

\section{RQ2: What Changes after Agent Adoption?}
\label{sec:rq2}

Author-level comparisons do not answer what happens to a repository after
agent adoption: a team may produce structurally similar changes but produce
more of them. We therefore ran a separate preregistered repository-level event
study. The registration froze the endpoints, exposure definitions, and
robustness battery before mining.

\subsection{Design}

The confirmatory panel contains 854 adopter repositories and 517 never-adopter
placebo controls, selected by the registered rules from the corpus as of
22~August~2026. The adoption event is the earliest of four observable
markers: an agent configuration file (\texttt{CLAUDE.md}), a
\texttt{copilot-swe-agent} pull request, a \texttt{codex/} branch, or a
\texttt{Co-Authored-By: Claude} trailer. Eligible repositories have at least
50 commits before adoption and at least 12 weeks of history on each side.
Placebo repositories are never-adopters; each receives a pseudo-adoption date
drawn deterministically from the adopter date distribution, subject to at
least 12 weeks of history on both sides, and passes through the identical
machinery. Placebo dates are calendar-matched, not activity-matched.

For each endpoint we estimate post/pre rate ratios under three exposure
definitions: per week (volume), per commit, and per KLOC added (per-unit).
Confidence intervals use owner-clustered bootstrap resampling; we additionally
evaluate leave-one-owner-out (LOO) stability, pre-trends on the pre-period,
and donut windows of $\pm4$ weeks around adoption. The primary endpoint, P1,
is the union of the per-commit drift events measured at weekly granularity:
complexity growth, boolean-field accretion, flag-argument growth, clone
introduction, SATD, and added include-graph edges. The rare gate-grade
structural events (new cycles, god headers) occur in well under 0.1\% of
commits and are too sparse for weekly rate estimation; they are analyzed at
commit level in Section~\ref{sec:rq1}.

\paragraph{Deviations from the registration.} The registration specified a
Poisson pseudo-maximum-likelihood (PPML) estimator with repository and
date-group$\times$event-week fixed effects (M2). Drift per commit has a
calendar-time shape in both arms: it rises through 2025, peaks in
January--March 2026, and falls afterwards. Placebo windows straddle that peak,
so the placebo post/pre ratio is below one (P1: 0.72) and an event-time
comparison of the two arms mixes adoption with calendar time. We therefore
report as the main estimate a PPML with repository and calendar-week fixed
effects (M3), which compares adopters with placebo repositories in the same
calendar week; the registered M2 estimate is given for P1 in the text. The
estimates submitted for review used pooled post/pre rate ratios on a smaller
panel (361 adopters, 265 placebo) and an earlier analyzer version; the panel
was rebuilt with the registered rules and replayed with archcheck~0.1.8.

Within each repository-week at most ten commits, chosen by hash, were
analyzed, and each is weighted by the number of commits in that week divided
by the number analyzed (210{,}594 analyzed commits). Excluded by outcome-blind
rules: 20{,}480 commits whose analysis exceeded the memory cap (9.7\%;
adopters 14.5\% before and 13.6\% after the event, placebo 9.9\% and 12.0\%),
1{,}656 with unreadable Git objects, 428 bulk imports, 75 pure file moves
($\ge$20 renamed C/C++ files and at least half of the changed C/C++ files),
and 360 for other technical reasons. A pure file move otherwise reports every
include edge of the moved files as new. KLOC counts added C/C++ lines.

The registered clone endpoint P4 is clone \emph{stock}; the panel measures
clone \emph{introductions}, and stock is measured only in the replication of
Section~\ref{sec:agarwal-repl}. TOST equivalence tests at $[0.80, 1.25]$ are
inconclusive for every per-commit and per-KLOC ratio. Not run: the
year-earlier and Dockerfile-adoption placebo controls and the pure-agent
stratum S1.

\begin{table*}[t]
\caption{Adoption event study (854 adopters, 517 placebo). Main estimate:
PPML incidence-rate ratio of the adopter post period, repository and
calendar-week fixed effects, owner-clustered 95\% CI. Adopter and placebo
columns: post/pre rate ratio per week of each arm. Pre-trend: slope of the
adopter pre period (SE) and $p$ of the joint pre-period test of the
registered event-study PPML.}
\label{tab:event}
\small
\begin{tabular}{@{}lllllll@{}}
\toprule
Endpoint & /week [95\% CI] & /commit [95\% CI] & /KLOC [95\% CI] & Adopter & Placebo & Pre-trend \\
\midrule
P1 union      & \textbf{1.93} [1.42, 2.62] & 1.24 [0.97, 1.58] & 1.14 [0.96, 1.35] & 1.67 & 0.72 & $+0.033$ (0.018), $p=0.87$ \\
P2 complexity & \textbf{1.79} [1.53, 2.09] & \textbf{1.17} [1.00, 1.37] & \textbf{1.14} [1.02, 1.27] & 1.79 & 0.86 & $+0.0035$ (0.0011), $p=0.06$ \\
P3 flag       & \textbf{1.41} [1.07, 1.85] & 0.93 [0.73, 1.19] & 0.92 [0.75, 1.14] & 1.32 & 0.73 & $+0.0015$ (0.0013), $p=0.004$ \\
P4 new clone  & \textbf{1.75} [1.37, 2.22] & 1.09 [0.89, 1.33] & 1.12 [0.92, 1.37] & 1.54 & 0.62 & $+0.017$ (0.004), $p=0.64$ \\
\bottomrule
\end{tabular}
\end{table*}

\subsection{Results}

Table~\ref{tab:event} summarizes the confirmatory results. After adoption the
weekly rate of P1 drift events is 1.93\x{} that of placebo repositories in the
same calendar weeks (95\% CI $[1.42, 2.62]$); the registered M2 estimator
gives 1.84\x{} ($[1.08, 3.13]$). Commit volume rises in step: census commits
in the window rise 1.39\x{} in adopters and fall to 0.94\x{} in placebo
repositories. Per commit, P1 rises 1.24\x{} ($[0.97, 1.58]$) and per KLOC
1.14\x{} ($[0.96, 1.35]$). P1 is heavy-tailed: added include edges and new
clones are most of its events, and single commits dominate some repositories;
winsorized at the 99th percentile of repository-week counts, the adopter/placebo
ratio of per-week rate ratios is 1.59.

Complexity (P2) is the least heavy-tailed endpoint and the one where the
per-unit rate also rises: 1.79\x{} per week, 1.17\x{} per commit
($[1.00, 1.37]$), 1.14\x{} per KLOC ($[1.02, 1.27]$). Its adopter pre-period
has a positive slope ($+0.0035$, SE $0.0011$), and the joint pre-period test
gives $p=0.06$, so part of the P2 increase may continue a trend that began
before adoption. In the registered dose--response split by post-adoption
AI-commit share, weekly complexity growth rises with share (1.35\x{} for zero,
1.43\x{} for low, 2.71\x{} for high), while per-KLOC rates do not (1.21\x{},
1.07\x{}, 1.22\x{}). Clone introductions (P4) and flag-argument growth (P3)
rise per week (1.75\x{}, 1.41\x{}) but not per commit (1.09\x{}, 0.93\x{}). The
pre-period of P3 fails the joint test ($p=0.004$).

\subsection{Independent replication on another study's repositories}
\label{sec:agarwal-repl}

As an external anchor, we replayed our full drift pipeline on the 21 C and C++
repositories from the dataset of Agarwal, He, and
Vasilescu~\cite{agarwal2026}, using their published per-repository adoption
dates: the same repositories, an independent instrument. Nineteen
repositories yield a usable pre/post split. The volume direction replicates:
median post/pre ratios are 1.56\x{} for commits per week and 2.03\x{} for
added KLOC per week. Weekly complexity events rise (median 1.49\x);
per-commit complexity does not (median 0.91\x), unlike the 1.17\x{} of the
panel. The registered clone \emph{stock} endpoint,
measured as whole-repository duplication snapshots at the window edges, is
null: median stock change $\approx$0\% ($n=19$; nearly every established
repository within $\pm$12\%). Rising clone \emph{introductions} with a flat
clone \emph{stock} are consistent: introductions are partially offset by
deletions and refactoring, and stock is dominated by pre-existing code. The
per-repository table is given in Appendix~\ref{app:agarwal}.

\medskip
\noindent\textbf{RQ2 answer.} Agent adoption is associated with a higher rate
of architectural drift per week, mostly through more commits. Per commit and
per KLOC the union endpoint rises less, with intervals that include one;
complexity rises per unit as well (1.17\x{} per commit), with a pre-period
trend that weakens this reading. The per-week direction is consistent on an
independent list of 19 repositories.

\section{RQ3: Review Capacity}
\label{sec:rq3}

A throughput-driven increase matters only if quality-control mechanisms scale
with the new volume. We measured human review before and after adoption with
the same adopter/placebo framing (owner-clustered bootstrap, LOO, donut,
pre-trends). Up to eight mainline commits per repository and period were
sampled and classified through their pull requests (19{,}978 commits, 844
adopter and 500 placebo owners). Merged pull requests per repository were
counted for the 12 weeks on each side of the event (1{,}372 repositories), and
for the 9{,}050 merged pull requests of the sample we collected every review,
inline comment, and conversation comment, excluding bots and the pull-request
author.

\begin{table*}[t]
\caption{Review capacity before/after adoption. Ratio: adopter post/pre over
placebo post/pre, owner-clustered bootstrap 95\% CI. Volume and comment rows
are exploratory (not registered).}
\label{tab:review}
\small
\begin{tabular}{@{}lllll@{}}
\toprule
Endpoint & Adopter pre $\to$ post & Placebo pre $\to$ post & Ratio [95\% CI] & $p$ \\
\midrule
R1: PRs w/o recorded human review (\%) & 38.7 $\to$ 42.5 & 43.1 $\to$ 44.6 & 1.06 [0.96, 1.18] & 0.29 \\
R2: direct pushes (\%)                  & 46.3 $\to$ 43.4 & 63.3 $\to$ 62.5 & 0.95 & --- \\
R3: PR size (median added lines)        & 43 $\to$ 60     & 41 $\to$ 39     & 1.47 [1.02, 2.14] & 0.04 \\
Merged PRs, 12 weeks                    & 60{,}231 $\to$ 78{,}170 & 17{,}950 $\to$ 16{,}885 & 1.38 [1.21, 1.58] & $<$0.001 \\
Human comments per PR                   & 2.79 $\to$ 2.62 & 2.36 $\to$ 3.05 & 0.73 [0.52, 0.99] & 0.05 \\
\bottomrule
\end{tabular}
\end{table*}

After adoption, adopters merge 1.30\x{} more pull requests while placebo
repositories merge 0.94\x{}, a ratio of 1.38 ($[1.21, 1.58]$). Pull requests
also become larger (median added lines 43 to 60; ratio 1.47). The number of
human review comments per pull request falls to 0.73\x{} of the placebo
trajectory ($[0.52, 0.99]$): the total number of human comments in adopting
repositories stays at the level of placebo repositories (ratio 1.00,
$[0.71, 1.44]$) while the number of pull requests rises. The share of pull
requests merged with no recorded human review rises from 38.7\% to 42.5\%;
against the placebo trajectory this is 1.06\x{} ($[0.96, 1.18]$), not
significant. Fewer changes bypass pull requests after adoption (direct pushes
46.3\% to 43.4\%), so the change happens inside the pull-request workflow.

The registered R5 endpoint, the share of drift-carrying commits merged without
human review, was not reproduced on the re-run panel (ratio 1.02,
$[0.91, 1.14]$); the submitted value (1.10) came from drift commits sampled
directly, and the re-run joins the review sample to the replay sample.

The decrease is not explained by reviewers spending their effort selectively on
agent code: within the same
repositories, agent-authored PRs receive no significantly different review
effort than human PRs on any measured axis (review count, review comments,
change requests), consistent with prior findings that agentic PRs are not
reviewed more~\cite{duma2026reviews}, and evidence against the hypothesis that
C++ reviewers, specifically, filter agent code more strictly.

\medskip
\noindent\textbf{RQ3 answer.} Human review does not scale with the
post-adoption increase in change volume: adopters merge 1.38\x{} more pull
requests relative to placebo, and human review comments per pull request fall
to 0.73\x{}. The share of pull requests with no recorded human review rises
slightly, without a significant difference from placebo.

\section{Discussion}
\label{sec:discussion}

The three analyses resolve an apparent contradiction. Pooled, AI-attributed
changes look worse on several metrics. Controlling for repository and change
size, the author-level penalty disappears --- in merged history and as
submitted. Measuring repositories over time, drift still accelerates after
adoption. These observations are compatible because they use different
denominators:

\begin{itemize}
\item quality per comparable change, by author: approximately flat (RQ1);
\item change volume per unit time: higher;
\item architectural-drift events per unit time: higher, with a smaller
  per-change rise in complexity;
\item human review comments per pull request: lower.
\end{itemize}

The main engineering risk of agent adoption, on this evidence, is a capacity
mismatch between generation and verification. A review process designed for
human-scale change volume becomes underprovisioned even when individual
AI-assisted changes are no worse than human changes of comparable size.

\paragraph{One contested per-unit signal.} A single ``agents worse'' result in
our broader investigation survived its checks and deserves an honest status
note: a per-structure analysis (git-blame attribution over structures touched
by $\ge$4 commits) found boolean-field accretion (flag fields accumulating
in existing structs) in 9.8\% of agent-attributed structs versus 6.0\% of
others, roughly 2.3\x{} after manually removing false positives from both
arms. The adoption-level proxy for the same phenomenon (P3) is null per commit
(0.93\x{}) and fails the pre-period test, and the
per-struct method has not been re-run against the adoption cohorts. The result
is therefore method-dependent: unconfirmed rather than refuted, and we present
it as an open question rather than a claim.

\paragraph{Implications for tooling.} The volume channel changes the role of
static architecture checks. A CI gate should not attempt to classify whether
code was written by a human or an agent; it can remain author-agnostic:
compare the proposed revision to the baseline and reject only precise
structural regressions. Deterministic checks scale with change volume in a way
human attention does not. We also examined whether contemporary AI review bots
already provide this structural coverage. Of the 648 corpus commits that grew
an include cycle or added a cross-area dependency, 336 went through a pull
request, and 43 of those (38 pull requests in 17 repositories) were reviewed
by an AI review bot. Two bot reviews named the structural problem; none
blocked the merge. The sample is small and
conditioned on known regressions, so it is not a general estimate of bot
review quality, but it illustrates that repository-graph properties require
graph-aware checks rather than diff-local reasoning.

\section{Threats to Validity}
\label{sec:threats}

\paragraph{AI attribution and the detectability ceiling.} Metadata-based
attribution has high precision but incomplete recall. Unmarked AI use (IDE
completion, stripped trailers) contaminates the human class and biases author
comparisons toward the null. Let $r$ be the true agent/human rate ratio, $a$
the share of agent-labelled commits that are human, and $h$ the share of
human-labelled commits that are AI-assisted. The observed ratio is
\[
  \hat r = \frac{(1-a)\,r + a}{(1-h) + h\,r},
\]
so with $a=0.2$, $h=0.3$ a true $r=2.0$ is observed as $1.8/1.3 \approx 1.38$.
Table~\ref{tab:seven} gives the power of each sign test: for complexity and
clones the test detects agents being worse in 60\% of non-tied repositories,
for grown cycles only in 78\%. RQ1 is therefore a failure to detect an agent
disadvantage under observable attribution, not evidence of equivalence.

\paragraph{Task mix.} Size matching does not control what a commit does. In a
separate sample of 50 adopting repositories, 79\% of post-adoption commits
without an AI marker were made by authors who mark other commits as
agent-assisted, and 43 of the 50 repositories record agent use only as a
trailer on a human-authored commit. In the two repositories we inspected by
hand, the unmarked commits were configuration, build and documentation
changes, and the marked commits carried the C++ feature work. Within a repository, the
agent/human contrast therefore partly compares kinds of work. The direction of
this bias is not known.

\paragraph{Observational design.} Neither the within-repository comparison nor
the event study randomly assigns AI use. Adoption may coincide with
organizational changes that affect development volume. The placebo design,
pre-trend checks, donut windows, and leave-one-owner-out analysis reduce but
do not eliminate this concern: a shock that triggers both adoption and a
volume increase (new funding, a new maintainer, a release push) would produce
the same pattern. Deviations from the registration are listed in
Section~\ref{sec:rq2}.

\paragraph{Measurement validity.} \tool{} measures physical and local
structural properties, not semantic architecture. Include cycles and
cross-area dependencies are meaningful in C++~\cite{lakos1996,
coreguidelines}, but they do not capture all forms of design degradation, and
some signals are intentionally heuristic. In manual audits, complexity
findings usually indicated genuinely complex functions but rarely removable
debt, and the test co-evolution signal had limited precision; we accordingly
avoid interpreting every finding as a defect, and no conclusion in this paper
rests on a single category. The token-level cognitive-complexity scorer
overestimates on some preprocessor-heavy constructs; it does so
deterministically and identically for both author classes.

\paragraph{Sparse structural events.} New cycles and god headers are rare at
commit level; the corresponding sign tests contain many ties and have limited
power. Those rows mean ``no detected disadvantage,'' not superiority.

\paragraph{Repository selection.} The corpus requires substantial recent
activity, which samples survivors, and some agent activity, which enriches it
for adopters (Section~\ref{sec:data-attribution}). Projects with enforced
review were excluded, so the review-coverage levels in RQ3 are higher than
in open source generally; the before/after change is less affected.
Born-agentic repositories are excluded by construction. The 344 repositories
of RQ1 have substantial activity from both classes, which selects for heavy
adopters. Results may not generalize to private industrial repositories or
other review workflows.

\paragraph{Scope.} The physical-design categories are specific to C and C++
include structure; the volume-versus-quality decomposition is not. The
categories cover physical structure and local code signals, not semantic
architecture such as layering or API design. Agent capabilities change
quickly, so the per-commit results describe agents of 2024--2026; the
measurement pipeline can be rerun on later data.

\paragraph{Review observability.} GitHub records review events, not whether a
human privately read a diff. ``No recorded human review'' is an observable
workflow measure, and prior work cautions that review-presence metrics
overstate human oversight~\cite{duma2026reviews}. Comment counts measure
written review, not reading; a reviewer may approve after reading without
comment. The comment and pull-request volume analyses were added after the
registration.

\section{Related Work}
\label{sec:related}

\paragraph{Impact of coding agents.} Agarwal, He, and
Vasilescu~\cite{agarwal2026} study agent adoption on the AIDev dataset with a
staggered difference-in-differences design and report large, front-loaded
velocity gains alongside persistent quality-metric increases (static-analysis
warnings and aggregate cognitive complexity, roughly $+18\%$ and $+39\%$).
Their quality outcomes are aggregated at repository-month level and are
therefore compatible with the volume mechanism we isolate: totals rise with
delivered lines. He et al.~\cite{he2026cursor} find, for Cursor adoption, a
large but transient velocity increase and a persistent rise in warnings and
complexity, again without a per-change decomposition. Liu et
al.~\cite{liu2026debt} mine over 300{,}000 verified AI-authored commits and
track the lifecycle of the code-smell, bug, and security issues they
introduce. Randomized and
within-developer designs report substantial individual productivity gains from
AI assistance~\cite{cui2024copilot}. Our contribution is complementary: we
measure structural \emph{deltas} per commit, explicitly separate per-change
quality from per-time volume, and extend the evidence to C++ physical design
--- include-graph cycles, god headers, cross-area dependencies --- which none
of the above measure.

\paragraph{Motivating concern.} Dente et al.~\cite{dente2026constraintdecay}
show experimentally that agent adherence collapses as structural constraints
accumulate (constraint decay), motivating the hypothesis that agent code
should drift more. Our field data suggest the per-unit effect, if present in
the wild, is below the detectability ceiling of metadata attribution, while
the operative risk materializes through volume.

\paragraph{Review of AI-generated changes.} Duma et
al.~\cite{duma2026reviews} find that most AI-generated PRs receive no review
and that human involvement is often agent steering rather than evaluation;
adoption studies likewise report agent commits being larger than human
ones~\cite{robbes2026agentic}. Our RQ3 adds the before/after,
placebo-controlled trajectory of review coverage within adopting C++
repositories, measuring both pull-request volume and human review comments.

\paragraph{Mining methodology.} Robbes et al.~\cite{robbes2026perils} catalog
pitfalls of agent-activity heuristics; Khosravani and
Mockus~\cite{khosravani2026census} validate multi-method agent detection and
quantify single-signal undercounting. Both inform our conservative
attribution and its stated limits. The consequences of clone introduction for
maintenance are classic~\cite{juergens2009}; our per-commit clone-introduction
measurements give that literature a drift-rate reading.

\section{Conclusion}

Across 1{,}188 C++ repositories and 484{,}500 analyzed commits, we find no
evidence that AI-attributed commits are more architecturally damaging than
human commits of comparable size within the same repository. The strongest
observable author difference is that AI-attributed commits are larger. After
agent adoption, repositories accumulate architectural-drift events 1.93\x{}
faster per week than placebo repositories in the same calendar weeks, mostly
through more commits; complexity also rises per commit (1.17\x{}). Adopters
merge 1.38\x{} more pull requests relative to placebo, while human review
comments per pull request fall to 0.73\x{}.

Coding agents increase the rate of change, structural regressions arrive
faster, and the amount of human review stays where it was. For teams adopting
coding agents, deterministic, author-agnostic structural checks in CI scale
verification with generation volume.

\section*{Data Availability}
\label{sec:data}

The artifact is deposited on Zenodo under DOI 10.5281/zenodo.24440566
(concept DOI 10.5281/zenodo.22057749) and
contains the analysis code for every study (per-commit pipeline, event-study
estimators, review-capacity estimators), the preregistration, the raw
estimator outputs from which the tables are drawn, the tables with a
per-number provenance manifest, and the analyzer. The re-run per-commit
records of the event study (210{,}594 commits), the cohort definitions, the
review panel, and the pull-request counts are part of the deposit. The
per-commit dataset behind RQ1 was lost in a machine reinstall; its analysis
code and raw outputs are in the deposit, and it can be regenerated from the
public repositories.

\bibliographystyle{ACM-Reference-Format}
\bibliography{refs}

\appendix

\section{Per-Repository Replication Table}
\label{app:agarwal}

Table~\ref{tab:agarwal} lists the per-repository post/pre ratios for the
replication of Section~\ref{sec:agarwal-repl}. The per-repository rows were
printed by the replay at the 94\%-complete point of the run; all aggregate
medians reconcile with the completed run except weekly complexity (1.48
vs.\ 1.49). Two repositories with degenerate pre/post windows are excluded
from ratio statistics; complexity ratios are undefined where one period has no
complexity events.

\begin{table*}[h]
\caption{Replication on the C/C++ repositories of~\cite{agarwal2026}:
post/pre ratios per week, and duplication-stock change.}
\label{tab:agarwal}
\small
\begin{tabular}{@{}lrrrr@{}}
\toprule
Repository & Commits/wk & KLOC/wk & Cplx/wk & Dup stock \\
\midrule
forntoh/LcdMenu            & 9.77  & 12.28 & 28.8 & $-8\%$ \\
joaquimorg/UV-KX           & 12.91 & 5.02  & 9.4  & $+71\%$ \\
AMICI-dev/AMICI            & 2.13  & 7.96  & 6.33 & $-9\%$ \\
microsoft/onnxruntime      & 1.18  & 0.99  & 2.17 & --- \\
M-CreativeLab/jsar-runtime & 1.56  & 3.28  & 1.91 & $+159\%$ \\
rsyslog/rsyslog            & 6.01  & 91.05 & 1.81 & $+4\%$ \\
wolfSSL/wolfssl            & 1.35  & 2.64  & 1.55 & $+3\%$ \\
Z3Prover/z3                & 0.90  & 0.51  & 1.41 & $-4\%$ \\
shader-slang/slang         & 1.24  & 0.86  & 1.03 & $-1\%$ \\
chrxh/alien                & 1.60  & 1.37  & 0.74 & $-11.6\%$ \\
microsoft/ebpf-for-windows & 0.87  & 0.90  & 0.66 & $+5\%$ \\
DeForestt/aflat            & 0.64  & 4.26  & 0.62 & $-5\%$ \\
microsoft/react-native-windows & 0.89 & 0.64 & 0.43 & --- \\
ohhmm/openmind             & 0.22  & 0.14  & 0.00 & $-6\%$ \\
SavenkovIgor/TermGraph     & 4.06  & 20.95 & ---  & $-6\%$ \\
PlantSimulationLab/Helios  & 1.72  & 1.44  & ---  & $-7\%$ \\
effekseer/Effekseer        & 3.33  & 2.03  & ---  & $+0\%$ \\
aboys-cb/NepTrainKit       & 0.80  & 0.27  & ---  & $+8\%$ \\
aantropov/sailor           & 4.06  & 17.78 & ---  & $+3\%$ \\
\bottomrule
\end{tabular}
\end{table*}

\end{document}
